\chapter{Výzkumná část}
\label{chap:research}

Tato kapitola představuje praktickou realizaci výzkumu. Nejprve je definován produkt Berrie a~provedena analýza konkurenčního prostředí na českém trhu. Následně jsou popsány průběh a~výstupy jednotlivých fází výzkumu.

% ============================================================================
\section{Definice produktu Berrie}
\label{sec:res_product}

Berrie představuje koncept mraženého produktu prémiové kvality — mražené ovoce obalené v~dvojité vrstvě čokolády. Ovoce je zpracováno technologií šokového mražení (\gls{iqf}), která probíhá při teplotách $-35$ až $-40\,^{\circ}\mathrm{C}$ za silného proudění vzduchu. Díky tomu vznikají pouze malé mikrokrystaly ledu a~buněčná struktura ovoce zůstává zachovaná, což zajišťuje lepší chuťové a~texturní vlastnosti ve srovnání s~běžným domácím mražením (Grover \& Negi, 2023; Oldenburg et al., 2025).

\subsection{Produktové portfolio}

Základní varianty produktu Berrie zahrnují:
\begin{itemize}
  \item borůvky v~bílé čokoládě,
  \item jahody v~hořké čokoládě,
  \item banán v~mléčné čokoládě,
  \item ovocný mix.
\end{itemize}

Do budoucna se počítá s~rozšířením portfolia o~veganské varianty a~další druhy ovoce.

\subsection{Klíčové vlastnosti}

\begin{table}[H]
\centering
\caption{Klíčové vlastnosti produktu Berrie}
\label{tab:product_features}
\begin{tabularx}{\textwidth}{lX}
\toprule
\textbf{Vlastnost} & \textbf{Popis} \\
\midrule
Surovina & České ovoce, případně v BIO kvalitě \\
Technologie & Šokové mražení (IQF) \\
Čokoláda & Dvojitá vrstva (kombinace typů čokolády) \\
Balení & 150g (s možností rozšíření) \\
Distribuce & Rohlík.cz, supermarkety \\
Cílová cena & Konkurenceschopná vůči Franui ($\sim$160 Kč/150g) \\
\bottomrule
\end{tabularx}
\end{table}

% ============================================================================
\section{Analýza konkurence}
\label{sec:res_competition}

\subsection{Franui — hlavní konkurent}

Hlavním konkurenčním produktem na českém trhu jsou mražené maliny obalené v~čokoládě značky Franui. Původ značky sahá do argentinské Patagonie, kde zakladatel Diego Fenoglio navázal na rodinnou tradici v~čokoládovém průmyslu. Nápad na spojení lokálních divokých malin s~italskou čokoládou vznikl v~roce 2013, přičemž klíčovým faktorem úspěchu se stala technologie šokového mražení, která umožnila masovou distribuci produktu (Franui, 2024).

Na evropský trh Franui expandovalo ve velkém měřítku v~roce 2021. Pro snazší logistiku byla v~roce 2020 otevřena výrobní továrna ve Valencii, odkud se produkt distribuuje na evropský trh. K~rychlé expanzi zásadně přispěl virální marketing na sociálních sítích TikTok a~Instagram, kde se produktová videa šířila organicky díky vizuální přitažlivosti produktu.

\begin{table}[H]
\centering
\caption{Srovnání Berrie a Franui}
\label{tab:comparison}
\begin{tabularx}{\textwidth}{lXX}
\toprule
\textbf{Kritérium} & \textbf{Berrie} & \textbf{Franui} \\
\midrule
Původ & Česká republika & Argentina / Španělsko \\
Ovoce & Borůvky, jahody, banán, mix & Maliny \\
Čokoláda & Dvojitá vrstva & Dvojitá vrstva (bílá + mléčná/hořká) \\
Klíčová výhoda & Lokální původ, BIO, více druhů & Zavedená značka, silný branding \\
Cenová pozice & Konkurenceschopná & $\sim$160 Kč / 150g \\
Distribuce & Rohlík.cz + supermarkety & Omezená dostupnost \\
\bottomrule
\end{tabularx}
\end{table}

\subsection{Mezera na českém trhu}

Z~marketingového pohledu je Franui tzv.\ cross-category produktem, spojujícím prvky zmrzliny, pralinky a~zdravější ovocné pochutiny. Na českém trhu v~současnosti chybí přímá lokální alternativa k~tomuto typu produktu. Zatímco český spotřebitel zná řadu tradičních mražených pochutin (zmrzlina Míša, nanuky Ledňáček apod.), segment mraženého ovoce v~prémiové čokoládě je obsazen výhradně importovaným produktem.

Hlavní konkurenční výhody značky Berrie vůči Franui:
\begin{itemize}
  \item využití lokálních surovin (české ovoce, případně \gls{bio}),
  \item širší portfolio druhů ovoce (nejen maliny),
  \item lepší dostupnost v~českých obchodních řetězcích,
  \item potenciálně konkurenceschopnější cena díky absenci importních nákladů.
\end{itemize}

% ============================================================================
\section{Realizace kvalitativního výzkumu}
\label{sec:res_qualitative}

\textit{Tato sekce bude doplněna po realizaci skupinových diskusí. Bude obsahovat detailní popis průběhu obou focus groups, přepisy klíčových výroků respondentů a~identifikované tematické vzorce.}

\subsection{Průběh skupinových diskusí}

% TODO: Doplnit po realizaci výzkumu

\subsection{Klíčové výroky respondentů}

% TODO: Doplnit po realizaci výzkumu

% ============================================================================
\section{Realizace kvantitativního výzkumu}
\label{sec:res_quantitative}

\textit{Tato sekce bude doplněna po provedení dotazníkového šetření. Bude obsahovat popis sběru dat, návratnost dotazníku a~demografický profil vzorku.}

% TODO: Doplnit po realizaci výzkumu

% ============================================================================
\section{Realizace senzorické analýzy}
\label{sec:res_sensory}

\textit{Tato sekce bude doplněna po provedení Central Location Testu. Bude obsahovat podmínky testování, hodnocení textury, chuťové rovnováhy a~celkového senzorického dojmu.}

% TODO: Doplnit po realizaci výzkumu
